\section{System Operation}\label{sec:system}
This section describes one mission in words; Section~\ref{sec:model} gives the model and notation.

\textbf{Setting.} The protected area is a lowland forest with gentle relief, such as Cat Tien National Park in southern
Vietnam (Fig.~\ref{fig:scenario}). The park operator owns a network of battery-powered devices: camera traps that take
pictures on motion, acoustic and wildfire sensors, and a few AI nodes, i.e., smart cameras or edge boxes placed in
clearings where a solar panel can power a neural accelerator. The devices share an IEEE 802.15.4 mesh at 2.4\,GHz. A
cluster head (CH) in each cluster of cells also has a LoRa radio that reaches the ranger station. The ranger station at
the forest edge hosts the base station (BS), a LoRa gateway, and the take-off sites of two UAVs: a fixed-wing VTOL for
delivery (team~1) and a rotary-wing UAV with a camera for verification (team~2).

\textbf{Step 1: phase 0, offline.} Long before any incident, the BS
\begin{enumerate}
\item tiles the area into hexagonal cells, elects a cell leader (CL) in every cell and a gateway pair between adjacent
cells, and groups the cells into clusters;
\item selects the AI cells and assigns every ordinary device to one AI node within a hop limit, so that the background
data of the devices (crops or embeddings of the persons they detect) stay within the traffic and memory budgets;
\item precomputes, for a set of plausible incident areas, the flight plan of team~1 and the ground transport that
complements it, and distributes the per-cell part of the plan through the mesh.
\end{enumerate}
The selection is refreshed when devices fail or are added; the plans are cheap to recompute.

\textbf{Step 2: incident.} A person is reported missing, with a last-seen point and a description. The BS builds the
query from reference photographs supplied by relatives, fixes the incident area around the last-seen point, picks or
recomputes the plan, and broadcasts a LoRa wake-up message of a few bytes with the mission identifier and the take-off
time. From the plan, every device in the incident area computes when it must be awake.

\textbf{Step 3: phase 1, delivery.} Team~1 takes off and flies the planned pass at constant altitude, broadcasting the
query packets in a fixed order. Every cell under the track keeps what its devices hear. As soon as a cell holds the
segment it is planned to forward, it sends it towards the AI cells over the cell-level routes of the plan. A repair
step fills the gaps left by lost packets or by a track that deviates from the plan, e.g., because of wind. When an AI
node holds the complete query, it matches the query against its stored detections and those of its assigned devices,
and sends a score of a few bytes through its CH to the BS over LoRa.

\textbf{Step 4: phase 2, verification.} The BS turns the reports into a posterior over the cells. Team~2 takes off,
flies to the cell with the best ratio of posterior to travel time, verifies it with its own camera, and continues with
the next cell; the order is recomputed when new reports arrive. The mission ends when team~2 confirms the person; the
coordinates are passed to the ground team.

Two remarks on scope. First, a single UAV serves phase~1; several UAVs would need a coordination mechanism that we leave
for future work. Second, the scenario is restricted to flat or gently rolling forest, where the UAV can hold a constant
altitude and the ground can be treated as a plane. Mountain parks add terrain shadowing and terrain-following flight
(Section~\ref{sec:model}), which we treat as an extension.
